\section{Past, Present, Future}

As we close in on the \textbf{1001} posts at which point Gornahoor will come to its natural end, it is time to focus more clearly on its goals. At the current pace, that end will come in about 10 to 12 months. There are other projects envisioned beyond that, which have been languishing. For example,

\begin{itemize}
\item There is much more work to do with the Medtarot discussion list.

\item The project to create graphical images for the various figures described in Guenon's Symbolism of the Cross is on hold. A few years ago, a mathematician from Notre Dame offered to help, but he has since moved on. Perhaps there is someone else who is familiar with ray tracing of mathematical modeling software.

\item There are various translations projects pending, mostly of obscure writers.

\item The Gnosis group, which is more important than any of the above, needs to expand.

\item Finally, of course, there are my own personal meditations which cannot risk being lost in a sea of busyness.
\end{itemize}


\paragraph{Ultimate Goal}


The main point of Gornahoor is to explore Tradition including its social structures, metaphysical principles, and esoteric teachings, with the ultimate aim being its possible restoration. We are not promoting any specific religious teachings, although we will use them to illustrate the manifestations of Tradition. We have been focusing on the Medieval tradition for the practical matter that it is closer to us in time, customs, language, etc., so presumably it will be easier to grasp for the modern mind seeking to understand.

Now the authors that primarily interest us --- \textbf{Rene Guenon}, \textbf{Julius Evola}, \textbf{Guido de Giorgio}, \textbf{Ananda Coomaraswamy} --- all agree that the European Middle Ages, along with its spiritual teachings, constituted a valid tradition. We have endeavored to explore those. Our aim is not to demonstrate their truth, since that is simply assumed. We have demonstrated that the Medieval tradition was in continuity with the preceding European traditions, although with a deepening understanding.

Properly understood, it provides a complete teaching, so there is no necessity to follow alternative, and alien, traditions. Nevertheless, the variety of religions, with their apparently incompatible claims to truth, has been a scandal to modern man. Science and rationalistic philosophies have not provided the basis for understanding the meaning of life. Various revolutionary movements have caused commotions without increasing human happiness.

At the right moment, then, Rene Guenon arrived with an intellectually and spiritually satisfying teaching. Commonality in symbols, rites, social organization, and metaphysical principles explain spiritual systems better than the alternatives from psychology and anthropology.

Nevertheless, we cannot stop there. It is one thing to announce the ``transcendent unity of religions''; specifity is a different matter. Hence, that principle dictated the following subgoals

\begin{itemize}


\item To distinguish between exoteric and esoteric. One way is to show how theological concepts can be recast as metaphysical concepts.

\item To show that the Medieval Church (or the Nordic-Roman tradition, in Evola's terms) formed a legitimate tradition

\item To highlight the differences between the Medieval and contemporary churches

\item To show the continuity of the Medieval church with the historically prior pagan traditions

\item To point out the homologies between the Medieval tradition and other traditions

\item To show that the Hermetic tradition preserved the essential elements of the medieval tradition
\end{itemize}


\paragraph{The Fall}


We can use the ``fall of man'' as an example of a homology between various traditions. Sacred texts have higher interpretations beyond the literal, historical, or scientific interpretation. The esoteric understanding stands at a different level. This does not mean it is the ``correct'' interpretation in opposition to the ``literal'' interpretation; viz., it does not assume the literal meaning is false. However, for our purposes, it is primarily the esoteric meaning that is of interest.

So, in this case, the esoteric interpretation of the Fall involves the knowledge of one's states of consciousness rather than the exoteric knowledge of the activities of the primal couple in the past. According to the Council of Trent, these are the lasting effects of the original fall, which persist to our own times:

\begin{itemize}
\item The intellect is darkened. It becomes difficult to grasp the true nature of the Man, God, and the World.

\item The will is weakened. It is beset by concupiscence and negative emotions which guide it instead of the higher intellect
\end{itemize}


This is a simpler way to put it. The fall of man involves man in:

\begin{itemize}


\item Ignorance

\item Desire
\end{itemize}


Put this way, this is the summation of Buddhism. The way of liberation begins with ``right view'' (the elimination of ignorance by seeing reality as it is) and ends with the end of excessive or disordered desire. Similarly, the Vedanta teaches that man is in a state of avidya (ignorance) and limited to the desire for pleasure, worldly success, or sense of duty. The ``antidote'' is vidya and desire refocused on liberation.

That is not the whole story, of course, but it is the basis for a common understanding.

\paragraph{Regeneration and the Mirror of God}


The twin peaks of scholastic metaphysics were Saint Thomas Aquinas and \textbf{Saint Bonaventura}.

In his mediations on the \textit{Soul's Journey to God}, Bonaventura uses the six wings of the seraph from Ezekiel as a springboard to a series of meditations on God, similar to the technique used by Valentin Tomberg on the imagery of the Tarot. He begins with the knowledge of God gleaned from the created world. As such, it is a nature mysticism, which many poets, etc., have experienced. Unfortunately, it is usually considered the endpoint rather than the first leg of a longer journey.

It ends with meditations of God as Being and as the Good, before finally resting in God. These meditations are probably far from the typical believer's understanding of God.

As far as I know, in the 8 centuries since Bonaventura's meditations were written, the Franciscan orders have never developed a retreat format based on them. He is not describing a psychological process, as some commentators believe. They may advocate some techniques from cognitive psychology or other schools. Stay away.

The soul's journey is also very unlike the ``faith journey'' that so many today want to dwell on. Bonaventure is objective, esoteric, and God-centered, whereas the faith journey is often maudlin, emotional, and self-centered.

Bonaventura's journey is to lead back to the primordial state prior to the fall. Called ``regeneration'', this is likewise the goal of Hermetism, not to mention Guenon's frequent discussions on the primordial state. Another topic mentioned by Bonaventura is the necessity to ``polish the mirror'' of the soul in order for the soul to better reflect God's image. Again, this is the absolute starting point for Tomberg's meditations. So we see the Hermetic tradition has preserved the Medieval esoteric tradition. However, they diverged at some point, even reaching enmity. That rift must be healed.

That is because, in our time, Bonaventura's ideas are often forgotten. For example, what pope or bishop mentions the desirability of restoring the primordial state? Furthermore, it is simply assumed today that every human is the image of God without making the effort of ``polishing the mirror''.

\paragraph{Simony}


Simon Magus tried to purchase the power of the Holy Spirit. Simony is a perpetual temptation and is considered a serious moral failing. Nevertheless, the sale of ``spiritual wisdom'' is a big business today. People prefer to pay for some teachings that they can repeat back like parrots. True teachings are priceless, since they require a change in being.


\flrightit{Posted on 2015-10-06 by Cologero}

\begin{center}* * *\end{center}

\begin{footnotesize}
\begin{sffamily}
\texttt{Erasmus Seabert on 2015-10-06 at 23:27 said:}

You're doing tremendous work, these posts are invaluable, thank you!

\hfill

\texttt{Aleksandar Gueric on 2015-10-07 at 05:46 said:}

As someone whose parents are both Catholic and Eastern Orthodox, I never felt animosity towards any of them, as it's usually a case with the Catholics or Orthodox church people. But I am baptized in the Orthodox church and I wonder are there any true Traditionalists among them?

I know about Serbian theologians like saint Iustin Popovic and saint Nikolai Velimirovic who were very antimodernist, but I'm not sure could they be considered Traditionalists or valuable for Traditionalist studies. In my opinion, Dimitrije Ljotic was a great example of an Orthodox author who could be valuable for some Traditionalists, although he had to focus much on the world politics and had to abandon more mystical persue. In my mind he is a good example of a ''bridge'' between Catholicism and Orthodoxy.

Also, I read recently a critique of Dimitrije Ljotic by certain Orthodox author (forgot the book title), and although Church Fathers agree that ``his heart was in a good place'' and he was a ``highly moral individual'', ``he'', author continues, ``seemed to have accepted certain Roman-Catholic ideas or practice of personal, mystic conection with Christ'' which is, according to the author, not very Orthodox friendly, where great importance is in Church meetings or ``sabors''. I didn't quite understand it, because isn't hesychasm a form of solitary, mystical practice of meditation and unity with the Divine? I was pretty sure that Orthodoxy has it's own share of more esoteric, solitary and mystical work.

I personally am not too fond of Church gatherings and crowds and find the practice of solitary, esoteric work much better, especially in these days. Also I have no problem reading saint Aquinas or other Catholic authors or authors who deal with Hermeticism, Oriental metaphysics etc. I basically grew up reading Franz Bardon, before finding Guenon or Evola. Just this would make me a heretic in the eyes or many Church goers. This is why I am now inwardly completely centered around Traditionalism as a whole, complete metaphysical system and care little about exoteric dogmatic separations of ''world religions''.

I will now try to find works of de Giorgio and saint Bonaventura, since I haven't read anything from them. Thanks for the recommendations!

\hfill

\texttt{Manwë on 2015-10-07 at 13:03 said:}

Philip Sherrard argues that catholicism deviates from the christian tradition. It would be interesting to read Gornahoor's take on Sherrard's criticism of St. Augustine and St. Thomas Aquinas.

From Theology to Philosophy in the Latin West

\url{http://www.myriobiblos.gr/texts/english/sherrard\_philosophy.html}


\hfill

\texttt{scardanelli on 2015-10-07 at 20:01 said:}

I'm glad to hear about the eventual resurrection of the medtarot discussion list. Let me know if I can help.

\hfill

\texttt{Mark Citadel on 2015-10-08 at 12:47 said:}

Will Gornahoor's articles remain available online or not? I am greatly saddened that this project will be discontinued.

\hfill

\texttt{Alistair Fraser on 2015-10-08 at 18:15 said:}

I thought this was apt and not aimed at you mark , just because your below it,

Its my tainted intuitive projection for mischievous reasons on all visitors to this site apart from the esteemed site host/hosts.

From Dirty Harry -Uh uh. I know what you're thinking. ``Did he fire six shots or only five?'' Well to tell you the truth in all this excitement I kinda lost track myself. But being this is a .44 Magnum, the most powerful handgun in the world and would blow your head clean off, you've gotta ask yourself one question: ``Do I feel lucky?'' Well, do ya, punk?

From Gornahoor Website -- Sniff Sniff, pours a glass of firewater, Uh Uh uh oh oh , yeah yeah yeah, cough cough , , i know what your all thinking , you've been coming here for years grabbing free information of the highest calibre inc speeches and insights and comments on the maestros of recent human traditional thinking.

Now to tell you the truth in all this search for the truth, i kinda lost track myself, I'm talking about somethings that have got to be said that are often left unsaid, I'm talking about what this sites effort costs to run, not at a profit , you understand , but just to run in and of itself, and I'm thinking , in here I've got the most powerful level of wisdom ever gathered in the traditional sense, and i have a love of that wisdom , its taken years to amass, i even have a cordial appreciation of some of you visitors , but I'm gonna ask you all directly just one time right here right now , a thing that is usually left unsaid
... ... .\\

Just tell me, no , show me... .how much DO U LOVE coming to this site and getting it all for nothing , at what point will it finally occur a clear thought in your minds, maybe i should DONATE a small snippet for general energy RUNNING COSTS

So now the code is broken , the unsaid has been said and now you punks have got to ask yourselfs one more question = Does your Love for this free platform feel Donatey , well does it ya punks hehehe

\hfill

\texttt{Cologero on 2015-10-08 at 18:33 said:}

Mr Citadel, the articles will remain online.

\hfill

\texttt{Cologero on 2015-10-08 at 18:43 said:}

@Manwe, Sherrard just embarrasses himself, since he chose the path of Strife. A better guide would be Vladimir Solovyov or Nicodemus who saw in Lorenzo Scupoli a spiritual guide, not a ``deviation from Christian tradition''.

That said, Thomism has some weak points; see, for example, the recent discussions on karma and dharma.

The question remains, then, about the ``cause'' of the modern world. I doubt it is presupposed in Thomsim or the filioque as Sherrard might think. Guenon and Evola had their own opinions. Mouravieff, I believe, is more on point. The exemplars of the Medieval period were the theologian and the knight. They towered over the ordinary in their respective domains, since they achieved self-mastery in the interior life.

However, the scientist and the enlightenment philosopher, the soldier with gunpowder, the merchants, came to dominate the exterior world. The theologian and the knight became irrelevant. The way out, then, will require those with a new power of integration.

\hfill

\texttt{Cologero on 2015-10-08 at 23:33 said:}

Thanks for the thought, Mr. Fraser. I've used a recent donation to purchase Evola's commentary on the Tao Te Ching.

However, what would be even better would be more blogs along the same lines; not necessarily some ideological purity, but more like a family dinner. Even comments made to other sites would be nice, beyond the small number of readers here.

Of course, there is the Gnosis study group which needs to expand. We are just finishing up a group for phase 1, so we will be looking to start up a new phase 1 group.

\hfill

\end{sffamily}
\end{footnotesize}

